\chapter{Complexity — Nondeterministic}
%%% generated by the blueprint pipeline from TCSlib.Complexity.TuringMachine.Nondeterministic
%%%%%%%%%%%%%%%%%%%%%%%%%%%%%%%%%%%%%%%%%%%%%%%%%%%%%%%%%%%%%%%%%%%%%%%%%%%%%%%
\section{Overview}

This module defines nondeterministic multi-tape Turing machines in the sense of
[AB09, \S 2.1.2]: machines with two transition functions $\delta_0$ and $\delta_1$, one
of which is selected by a choice bit at every step, so that a finite run is governed by a
choice word with one bit per step.  It provides the run semantics under a choice word, the
predicate that every branch halts within a given number of steps (monotone in the bound),
the bundled layer of machines with a finite state set, and the view of a deterministic
machine as a nondeterministic one that ignores its choices.

%%%%%%%%%%%%%%%%%%%%%%%%%%%%%%%%%%%%%%%%%%%%%%%%%%%%%%%%%%%%%%%%%%%%%%%%%%%%%%%
\section{Declarations}

\begin{definition}[Nondeterministic multi-tape Turing machine]
\lean{Turing.NDTM}
\label{Turing.NDTM}
\leanok
\uses{Turing.Action}
A nondeterministic Turing machine with a read-only input tape, $k$ work tapes, and a
write-only output tape, over an alphabet of symbols extended by a blank, consists of an
initial state $q_0$ and two transition functions $\delta_0$ and $\delta_1$, one for each
value of a choice bit.  Each $\delta_b$ maps the current state, the (possibly blank) symbol
under the input head, and the $k$ symbols under the work-tape heads to a one-step action:
a movement of the input head, a write-and-move instruction for each work tape, an optional
symbol to append to the output, and a successor state or the decision to halt.  Neither the
alphabet nor the state set is required to be finite.
\end{definition}

\begin{definition}[One step under a choice bit]
\lean{Turing.NDTM.stepWith}
\label{Turing.NDTM.stepWith}
\leanok
\uses{Turing.Action.apply, Turing.Cfg, Turing.Cfg.inputSymbol, Turing.Cfg.workTapeSymbols}
Let $M$ be a nondeterministic multi-tape Turing machine with transition functions
$\delta_0, \delta_1$, let $b \in \{0,1\}$ be a choice bit, and let $c$ be a configuration
of $M$ on a fixed input.  The step of $M$ from $c$ under $b$ returns $c$ unchanged if $c$
is halted (its state component is empty), whatever the value of $b$; otherwise it applies
to $c$ the action that $\delta_b$ prescribes for the current state, the symbol under the
input head, and the symbols under the $k$ work-tape heads.
\end{definition}

\begin{definition}[Initial configuration of a nondeterministic machine]
\lean{Turing.NDTM.initCfg}
\label{Turing.NDTM.initCfg}
\leanok
\uses{Turing.Cfg, Turing.Cfg.init}
The starting configuration of a nondeterministic multi-tape Turing machine $M$ on an input
word $x$: the machine is in its initial state $q_0$, the input head scans the first symbol
of $x$, all work tapes are blank with their heads at cell $0$, and the output tape is
empty.
\end{definition}

\begin{definition}[Run under a choice word]
\lean{Turing.NDTM.runWith}
\label{Turing.NDTM.runWith}
\leanok
\uses{Turing.Cfg, Turing.NDTM.stepWith}
Let $M$ be a nondeterministic multi-tape Turing machine and $c$ a configuration of $M$ on a
fixed input.  For a choice word $w = b_1 b_2 \cdots b_m \in \{0,1\}^*$, the run of $M$ from
$c$ under $w$ is the configuration obtained from $c$ by performing $m$ steps of $M$, the
$i$-th of them under the choice bit $b_i$.  Recursively, the run under the empty word is $c$,
and the run under $b w$ is the run under $w$ from the configuration reached by one step
of $M$ from $c$ under $b$.
\end{definition}

\begin{lemma}[Run under the empty choice word]
\lean{Turing.NDTM.runWith_nil}
\label{Turing.NDTM.runWith_nil}
\leanok
\uses{Turing.Cfg, Turing.NDTM.runWith}
\difficulty{1}
Let $M$ be a nondeterministic multi-tape Turing machine and $c$ a configuration of $M$ on
some input.  Running $M$ from $c$ under the empty choice word yields $c$ itself.
\end{lemma}

\begin{lemma}[Run under a choice word with a leading bit]
\lean{Turing.NDTM.runWith_cons}
\label{Turing.NDTM.runWith_cons}
\leanok
\uses{Turing.Cfg, Turing.NDTM.runWith, Turing.NDTM.stepWith}
\difficulty{1}
Let $M$ be a nondeterministic multi-tape Turing machine, $c$ a configuration of $M$ on some
input, $b \in \{0,1\}$ a choice bit, and $w$ a choice word.  Running $M$ from $c$ under the
word $b w$ gives the same configuration as running $M$ under $w$ from the configuration
reached by one step of $M$ from $c$ under the choice bit $b$.
\end{lemma}

\begin{lemma}[Run under a concatenated choice word]
\lean{Turing.NDTM.runWith_append}
\label{Turing.NDTM.runWith_append}
\leanok
\uses{Turing.Cfg, Turing.NDTM.runWith, Turing.NDTM.runWith_cons}
\difficulty{4}
Let $M$ be a nondeterministic multi-tape Turing machine, $c$ a configuration of $M$ on some
input, and $w, w'$ choice words.  Running $M$ from $c$ under the concatenation $w w'$ gives
the same configuration as first running $M$ from $c$ under $w$ and then running $M$ under
$w'$ from the configuration so reached.
\end{lemma}

\begin{lemma}[Halted configurations are fixed under every choice]
\lean{Turing.NDTM.stepWith_of_halt}
\label{Turing.NDTM.stepWith_of_halt}
\leanok
\uses{Turing.Cfg, Turing.NDTM.stepWith}
\difficulty{1}
Let $M$ be a nondeterministic multi-tape Turing machine, $b \in \{0,1\}$ a choice bit, and
$c$ a configuration of $M$ on some input whose state component is empty (the machine has
halted).  Then one step of $M$ from $c$ under $b$ leaves $c$ unchanged.
\end{lemma}

\begin{lemma}[Runs from a halted configuration stay put]
\lean{Turing.NDTM.runWith_of_halt}
\label{Turing.NDTM.runWith_of_halt}
\leanok
\uses{Turing.Cfg, Turing.NDTM.runWith, Turing.NDTM.runWith_cons, Turing.NDTM.stepWith_of_halt}
\difficulty{4}
Let $M$ be a nondeterministic multi-tape Turing machine and $c$ a configuration of $M$ on
some input whose state component is empty (the machine has halted).  Then for every choice
word $w$, running $M$ from $c$ under $w$ yields $c$ itself.
\end{lemma}

\begin{definition}[Halting within $t$ steps on every branch]
\lean{Turing.NDTM.HaltsWithin}
\label{Turing.NDTM.HaltsWithin}
\leanok
\uses{Turing.NDTM, Turing.NDTM.initCfg, Turing.NDTM.runWith}
Let $M$ be a nondeterministic multi-tape Turing machine, $x$ an input word, and
$t \in \mathbb{N}$.  The machine $M$ halts on $x$ within $t$ steps along every branch if,
for every choice word $w \in \{0,1\}^t$ of length exactly $t$, the configuration reached by
running $M$ under $w$ from its initial configuration on $x$ is halted, that is, has an
empty state component.
\end{definition}

\begin{theorem}[All-branch halting is monotone in the time bound]
\lean{Turing.NDTM.HaltsWithin.mono}
\label{Turing.NDTM.HaltsWithin.mono}
\leanok
\uses{Turing.NDTM, Turing.NDTM.HaltsWithin, Turing.NDTM.initCfg, Turing.NDTM.runWith_append, Turing.NDTM.runWith_of_halt}
\difficulty{4}
Let $M$ be a nondeterministic multi-tape Turing machine, $x$ an input word, and
$t \le t'$ natural numbers.  If for every choice word of length $t$ the run of $M$ under
it from the initial configuration on $x$ ends in a halted configuration, then the same
holds for every choice word of length $t'$.
\end{theorem}

\begin{definition}[Bundled finite nondeterministic Turing machine]
\lean{Turing.FinNDTM}
\label{Turing.FinNDTM}
\leanok
\uses{Turing.NDTM}
For an alphabet $\Sigma$, a finite nondeterministic Turing machine consists of a number
$k \in \mathbb{N}$ of work tapes, a finite state set with decidable equality, and a
nondeterministic multi-tape Turing machine with $k$ work tapes and those states, whose tape
cells hold either a symbol of $\Sigma$ or a blank.
\end{definition}

\begin{definition}[Deterministic machine as a nondeterministic one]
\lean{Turing.MultiTapeTM.toNDTM}
\label{Turing.MultiTapeTM.toNDTM}
\leanok
\uses{Turing.MultiTapeTM, Turing.NDTM}
Let $M$ be a deterministic multi-tape Turing machine with $k$ work tapes, initial state
$q_0$, and transition function $\delta$.  The associated nondeterministic machine has the
same tapes, states, and initial state $q_0$, and both of its transition functions
$\delta_0$ and $\delta_1$ are equal to $\delta$, so that its behaviour does not depend on
the choice bit.
\end{definition}

\begin{lemma}[Deterministic machine viewed nondeterministically: same start]
\lean{Turing.MultiTapeTM.toNDTM_initCfg}
\label{Turing.MultiTapeTM.toNDTM_initCfg}
\leanok
\uses{Turing.MultiTapeTM, Turing.MultiTapeTM.initCfg, Turing.MultiTapeTM.toNDTM, Turing.NDTM.initCfg}
\difficulty{1}
Let $M$ be a deterministic multi-tape Turing machine with transition function $\delta$,
let $N$ be the nondeterministic machine with the same initial state whose two transition
functions both equal $\delta$, and let $x$ be an input word.  Then the initial configuration
of $N$ on $x$ coincides with the initial configuration of $M$ on $x$.
\end{lemma}

\begin{theorem}[Deterministic machines ignore their choices]
\lean{Turing.MultiTapeTM.toNDTM_runWith}
\label{Turing.MultiTapeTM.toNDTM_runWith}
\leanok
\uses{Turing.Cfg, Turing.MultiTapeTM, Turing.MultiTapeTM.runFrom, Turing.MultiTapeTM.runFrom_succ_eq_step, Turing.MultiTapeTM.step, Turing.MultiTapeTM.toNDTM, Turing.NDTM.runWith, Turing.NDTM.runWith_cons}
\difficulty{4}
Let $M$ be a deterministic multi-tape Turing machine with transition function $\delta$, and
let $N$ be the nondeterministic machine with the same initial state whose two transition
functions both equal $\delta$.  For every configuration $c$ on a fixed input and every
choice word $w \in \{0,1\}^*$, the run of $N$ from $c$ under $w$ equals the configuration
reached by running $M$ from $c$ for $|w|$ steps.
\end{theorem}

\begin{definition}[Bundled deterministic machine as a nondeterministic one]
\lean{Turing.FinTM.toFinNDTM}
\label{Turing.FinTM.toFinNDTM}
\leanok
\uses{Turing.FinNDTM, Turing.FinTM, Turing.MultiTapeTM.toNDTM}
Let $\Sigma$ be an alphabet and $M$ a finite deterministic Turing machine over $\Sigma$ with
$k$ work tapes, a finite state set, and transition function $\delta$.  The associated finite
nondeterministic Turing machine has the same number $k$ of work tapes and the same state
set, and as underlying machine the nondeterministic machine with the same initial state
whose two transition functions both equal $\delta$.
\end{definition}
